\documentclass[10pt,a4paper]{amsart}
\usepackage[margin=19mm]{geometry}
\usepackage{amsmath,amssymb,mathtools}
\usepackage[scr=rsfs,cal=euler]{mathalfa}
\usepackage{xfrac}
\usepackage[unicode,hidelinks,bookmarks=false]{hyperref}
\newcommand{\ie}{i.e.\ }
\newcommand{\cf}{cf.\ }
\newcommand{\resp}{respectively\ }
\newcommand{\Subsection}{\subsection}
\providecommand{\nobreakdash}{\nobreak\hbox{-}}
\setlength{\emergencystretch}{2em}
\multlinegap0pt
\usepackage{mathtools}
\usepackage{bm}
\usepackage{algorithm}\usepackage{algpseudocode}
\usepackage{amssymb}
\usepackage{tikz}
\usepackage[normalem]{ulem}
\newtheorem{proposition}{Proposition}
\title{The algebraic structures of social organizations: \\ the operad of cooperative games}
\author{Dylan Laplace Mermoud, Victor Roca i Lucio}
\date{Source: arXiv:2507.01969v2 (CC BY 4.0)}
\begin{document}
\maketitle

\noindent\emph{Source and licence.} The algebraic structures of social organizations: \\ the operad of cooperative games. Version arXiv:2507.01969v2 (CC BY 4.0), distributed by arXiv under the Creative Commons Attribution 4.0 International licence, http://creativecommons.org/licenses/by/4.0/, as recorded in the arXiv licence metadata for this exact version and shown on the abstract page https://arxiv.org/abs/2507.01969v2. Full licence notice: \texttt{licenses/arxiv-2507.01969-CC-BY-4.0.txt}.\par\smallskip

\noindent\textit{Editorial scope.} Counted unit: the article's proposition surjective (source0.tex lines 2381--2387). The article's complete original proof of it is delivered with it (source0.tex lines 2389--2456, one \texttt{proof} environment of 68 lines). No mathematics is written, repaired or re-derived: the statement and the proof are the article's own lines, in the article's own notation, and no retained line is substituted or re-flowed.  No further source-local context is needed: every cross-reference of every retained line resolves inside the two delivered spans.  Nothing was omitted from any retained span, so the package carries no omission record.  No retained line contains a citation, so the wrapper carries no bibliography and no bibliography entry.  No retained line contains a figure environment, an image-inclusion command or drawing code, so no image is delivered and the package declares no asset dependency.  Editorial formatting only, in addition: an AMS \texttt{amsart} preamble that restates the article's own declarations and macros used by the retained lines, namely the environments \texttt{proposition} on separate counters, together with the standard packages those lines need.  The retained environments print in this document as Proposition 1 in order of appearance, whereas the article prints the same items with the numbering of its own full document.  The complete native source of the article as distributed in the arXiv e-print for this exact version is bundled unchanged as \texttt{sources/180917/source0.tex}.
\begin{proposition}\label{prop: tensor of imputations is surjective}
Let \(\Gamma_A = (A, \alpha)\) and \(\Gamma_B =(B, \beta)\) be two non-negative games, and let \(i \in A\). If $\alpha(\{i\}) = 0$ and \(\beta(\{l\}) = 0\) for all \(l \in B\), the map
\[
(- \otimes_i -): I(\Gamma_A) \times I(\Gamma_B) \longrightarrow I(\Gamma_A \circ_i \Gamma_B)~, 
\]
is surjective. 
\end{proposition}

\begin{proof}
The strategy of the proof is to first show that all the vertices of $I(\Gamma_A \circ_i \Gamma_B)$ can be reached by the partial tensor construction of the vertices of \(I(\Gamma_A)\) and \(I(\Gamma_B)\) and then deduce the general case. The vertices of \(I(\Gamma_A)\) and \(I(\Gamma_B)\) are respectively given by
    \[
    a^{(p)} \coloneqq \alpha_{\text{sp}} \mathbf{1}^{\{p\}} + \sum_{k \in A \setminus \{i\}} \alpha(\{k\}) \mathbf{1}^{\{k\}} \qquad \text{and} \qquad b^{(q)} \coloneqq \beta_{\text{sp}} \mathbf{1}^{\{q\}} = \beta(B) \mathbf{1}^{\{q\}},
    \]
since, by assumption,  \(\beta(\{l\}) = 0\) for all \(l \in B\). We first assume that \(p \neq i\). So, because \(a^{(p)}_i = 0\), the partial tensor product gives
    \[
    a^{(p)} \otimes_i b^{(q)} = \beta(B) a^{(p)}. 
    \]
    On the other hand, whenever \(p \in A \setminus \{i\}\), we have 
    \[
    c^{(p)} = \gamma_{\text{sp}} \mathbf{1}^{\{p\}} + \sum_{m \in A \diamond_i B} \gamma_m \mathbf{1}^{\{m\}} = \gamma_{\text{sp}} \mathbf{1}^{\{p\}} + \sum_{k \in A \setminus \{i\}} \gamma(\{k\}) \mathbf{1}^{\{k\}} + \sum_{l \in B} \gamma(\{l\}) \mathbf{1}^{\{l\}}. 
    \]
    Let us evaluate the coefficients in this sum. If \(m \in A \setminus \{i\}\), we have \(\gamma(\{m\}) = \beta(B) \alpha(\{m\})\), and if \(m \in B\), we have \(\gamma(\{m\}) = \alpha(\{i\}) \beta(\{m\}) = 0\). From this, we compute 
    \[ \begin{aligned}
    \gamma_{\text{sp}} = \gamma_C - \sum_{m \in A \diamond_i B} \gamma_m & = \alpha(A) \beta(B) - \sum_{k \in A\setminus \{i\}} \gamma(\{k\}) - \sum_{l \in B} \gamma(\{l\}) \\
    & = \alpha(A) \beta(B) - \beta(B) \sum_{k \in A \setminus \{i\}} \alpha(\{k\}) - \alpha(\{i\}) \sum_{l \in B} \beta(\{l\}) \\
    & = \beta(B) \alpha_{\text{sp}}. 
    \end{aligned} \]
    Then, 
    \[
    c^{(p)} = \beta(B) \alpha_{\text{sp}} \mathbf{1}^{\{p\}} + \beta(B) \sum_{k \in A \setminus \{i\}} \alpha(\{k\}) \mathbf{1}^{\{k\}} = \beta(B) a^{(p)} = a^{(p)} \otimes_i b^{(q)}. 
    \]
    Notice that in this case, \(c^{(p)}\) is the partial tensor product between \(a^{(p)}\) and any imputation of \(I(\Gamma_B)\). We finish the first part of the proof by looking at the tensor product
    \[ \begin{aligned} 
    a^{(i)} \otimes_i b^{(q)} & = \alpha_{\text{sp}} \mathbf{1}^{\{i\}} \otimes_i b^{(q)} + \sum_{k \in A \setminus \{i\}} \alpha(\{k\}) \mathbf{1}^{\{k\}} \otimes_i b^{(q)} \\
    & = \alpha_{\text{sp}} \beta(B) \mathbf{1}^{\{q\}} + \beta(B) \sum_{k \in A \setminus \{i\}} \alpha(\{k\})
    \mathbf{1}^{\{k\}}. 
    \end{aligned} \]
    On the other hand, for \(q \in B\), we have 
    \[
    c^{(q)} = \beta(B) \alpha_{\text{sp}} \mathbf{1}^{\{q\}} + \beta(B) \sum_{k \in A \setminus \{i\}} \alpha(\{k\}) \mathbf{1}^{\{k\}} = a^{(i)} \otimes_i b^{(q)}, 
    \]
    and so any vertex of \(I(\Gamma_A \circ_i \Gamma_B)\) is a partial tensor product of vertices of \(I(\Gamma_A)\) and \(I(\Gamma_B)\). 

    \medskip

    By convexity, any imputation \(z \in I(\Gamma_A \circ_i \Gamma_B)\) can be written as a convex combination
    \[
    z = \sum_{m \in A \diamond_i B} \lambda_m c^{(m)} = \sum_{k \in A \setminus \{i\}} \lambda_k c^{(k)} + \sum_{l \in B} \lambda_l c^{(l)}, 
    \]
    with \(\lambda_m, \lambda_k, \lambda_l\) all non-negative and \(\sum_{m \in A \diamond_i B} \lambda_m = \sum_{k \in A \setminus \{i\}} \lambda_k + \sum_{l \in B} \lambda_l = 1\). Set \(\lambda_A = \sum_{k \in A \setminus \{i\}} \lambda_k\) and \(\lambda_B = \sum_{l \in B} \lambda_l\). Define, for all \(k \in A \setminus \{i\}\) and \(l \in B\), the new coefficients \(\lambda'_k = \lambda_A^{-1} \lambda_k\) and \(\lambda'_l = \lambda_B^{-1} \lambda_l\). Hence, 
    \[
    \sum_{k \in A \setminus \{i\}} \lambda'_k = \lambda_A^{-1} \sum_{k \in A \setminus \{i\}} \lambda_k = 1, \qquad \text{and} \qquad \sum_{l \in B} \lambda'_l = \lambda_B^{-1} \sum_{l \in B} \lambda = 1. 
    \]
    We can finally write \(z\) as 
    \[
    z = \lambda_A \sum_{k \in A \setminus \{i\}} \lambda'_k c^{(k)} + \lambda_B \sum_{l \in B} \lambda'_l c^{(l)}, \qquad \text{with} \quad \lambda_A + \lambda_B = 1, \quad \lambda_A, \lambda_B \geq 0. 
    \]
    Using the fact that the vertices can split into partial tensor products, we have 
    \[ \begin{aligned} 
    z & = \sum_{k \in A \setminus \{i\}} \lambda_k a^{(k)} \otimes_i b^{(q_k)} + \sum_{l \in B} \lambda_l a^{(i)} \otimes_i b^{(l)} \\
    & = \sum_{k \in A \setminus \{i\}} \lambda_k a^{(k)} \otimes_i b^{(q_k)} + \lambda_B a^{(i)} \otimes_i \sum_{l \in B} \lambda'_l b^{(l)}. 
    \end{aligned} \]
    For \(k \in A \setminus \{i\}\), the vertex \(c^{(k)}\) decomposes using \(a^{(k)}\) and any imputation we want from \(I(\Gamma_B)\). Because \(\sum_{l \in B}  \lambda'_l = 1\) and \(\lambda'_l \geq 0\) for all \(l \in B\), we have that \(y = \sum_{l \in B} \lambda'_l b^{(l)}\) is a convex combination of the vertices of \(I(\Gamma_B)\), and by convexity, is itself a imputation, i.e., \(y \in I(\Gamma_B)\). Thus, 
    \[
    z = \sum_{k \in A \setminus \{i\}} \lambda_k a^{(k)} \otimes_i y + \lambda_B a^{(i)} \otimes_i y = \left( \lambda_B a^{(i)} +  \sum_{k \in A \setminus \{i\}} \lambda_k a^{(k)} \right) \otimes_i y. 
    \]
    Define \(x = \lambda_B a^{(i)} + \sum_{k \in A \setminus \{i\}} \lambda_k a^{(k)}\). By definition of \(\lambda_B\), we have that 
    \[
    \lambda_B + \sum_{k \in A \setminus \{i\}} \lambda_k = \lambda_B + \lambda_A = 1. 
    \]
    Moreover, \(\lambda_k \geq 0\) for all \(k \in A \setminus \{i\}\). So, \(x\) is a convex combination of the vertices of \(I(\Gamma_A)\), and therefore \(x \in I(\Gamma_A)\) by convexity. So, we have that 
    \[
    z = x \otimes_i y, \qquad \text{with} \qquad x \in I(\Gamma_A) \quad \text{and} \quad y \in I(\Gamma_B), 
    \]
    which concludes the proof. 
\end{proof}

\end{document}
